\section{Introduction}
\label{sec:intro}

Removing transformer blocks and then continuing pretraining (CPT, often called
``healing'') is an attractive route to smaller language models
~\citep{gromov2024unreasonable,shortenedllama,men2024shortgpt,
yang2024laco}. Held-out perplexity is a natural training signal, but using its
improvement as evidence of target recovery makes a stronger measurement
assumption: that likelihood and the target evaluation recover together after
the structural intervention.

Compression research has already established that perplexity can miss changes
in parametric knowledge, downstream behavior, and safety
~\citep{costcompression,jaiswal2024truth,beyondperplexity}. Depth-pruning work
also reports recovery curves and loss--task gaps
~\citep{gromov2024unreasonable,shortenedllama,minitron,iterabre}. The open
measurement problem is therefore not whether the metrics can disagree, but
\emph{how to audit a recovery claim}: is the apparent gap a property of the
training path, the evaluation interface, or the exact pruned construction?

We answer this question with a densely instrumented OLMo-2-1124-7B case study
~\citep{olmo2}. Among the post-depth-pruning studies we identified, none
jointly combines (i) same-source held-out perplexity and within-path
checkpoints, (ii) answer-letter and complete-option MMLU on the same 14,042
items, (iii) three no-retrieval closed-book QA datasets, and (iv) an intact CPT
branch, same-shape null points, and an alternative 16-layer construction.
This combination turns a familiar ``PPL versus task'' comparison into a
multi-axis test of proxy, interface, and construction.

\begin{figure*}[t]
\centering
\includegraphics[width=\textwidth]{figures/fig1_pruneheal_autofigure.pdf}
\caption{\textbf{A multi-axis audit changes the recovery verdict.}
\textbf{Left:} keep14 improves from 128k to 200k, yet the large intact-base
MMLU gap remains. \textbf{Right:} the random model exposes a high
complete-option floor, whereas ShortGPT reaches a much stronger endpoint at the
same nominal depth and displayed 200k-step budget; full32 supplies the available 25k
intact-CPT reference. PPL uses a same-source held-out Dolmino shard. Each
trained construction is one run and token counts are nominal presentations.
The exact unmatched factors are reported in Table~\ref{tab:main-results} and
Limitations.}
\label{fig:proxy-validity}
\end{figure*}

Figure~\ref{fig:proxy-validity} previews three measurement findings. First,
keep14 continues to improve late in training without approaching the intact
target scores. Second, complete-option MMLU makes the compressed models look
stronger, but the same interface also gives a random model a high floor. Third,
at the same nominal depth and displayed steps, ShortGPT exceeds keep14 by
15.5 points on letter MMLU but only 1.8 points on content MMLU while changing
four construction factors together. Construction and interface are therefore
both essential parts of the recovery result.

\paragraph{Study design.}
The principal construction retains the first 14 of 32 pretrained blocks,
appends two fresh blocks, and trains for 200k steps. We compare it with the
intact base, a full32 CPT checkpoint at 25k, frozen and random same-shape
operating points, and a non-contiguous ShortGPT-16 model at 200k. All PPL is
in-domain on a disjoint shard from the same Dolmino/DCLM source. Each trained
construction is one run; full32 has no 200k checkpoint, and ShortGPT changes
four structural choices together. We state these boundaries once here and
concentrate the remaining qualifications in Limitations.

\paragraph{Contributions.}
\begin{itemize}[leftmargin=1.2em,itemsep=2pt]
\item \textbf{A multi-interface recovery audit.} We jointly track same-source
and out-of-domain likelihood, paired letter/content MMLU, three closed-book QA
datasets, and continuation-stream overlap, with item-level uncertainty for the
load-bearing MMLU comparisons.
\item \textbf{Null and construction diagnostics.} An intact CPT reference,
random and frozen operating points, and ShortGPT separate three questions that
a single loss curve conflates: training-domain fit, interface validity, and
construction dependence.
\item \textbf{Four actionable findings and an audit protocol.} We quantify the
late keep14 target gap, the random content-score floor, the
ShortGPT--keep14 gaps of 15.5 letter points and 1.8 content points at the same
nominal depth/displayed steps, its consistent closed-book advantage, and the
much smaller short-horizon full32 deficits;
we turn them into a six-axis reporting checklist for future recovery studies.
\end{itemize}
